% Options for packages loaded elsewhere
\PassOptionsToPackage{unicode}{hyperref}
\PassOptionsToPackage{hyphens}{url}
\documentclass[
  10pt,
]{article}
\usepackage{xcolor}
\usepackage[margin=1in]{geometry}
\usepackage{amsmath,amssymb}
\setcounter{secnumdepth}{-\maxdimen} % remove section numbering
\usepackage{iftex}
\ifPDFTeX
  \usepackage[T1]{fontenc}
  \usepackage[utf8]{inputenc}
  \usepackage{textcomp} % provide euro and other symbols
\else % if luatex or xetex
  \usepackage{unicode-math} % this also loads fontspec
  \defaultfontfeatures{Scale=MatchLowercase}
  \defaultfontfeatures[\rmfamily]{Ligatures=TeX,Scale=1}
\fi
\usepackage{lmodern}
\ifPDFTeX\else
  % xetex/luatex font selection
\fi
% Use upquote if available, for straight quotes in verbatim environments
\IfFileExists{upquote.sty}{\usepackage{upquote}}{}
\IfFileExists{microtype.sty}{% use microtype if available
  \usepackage[]{microtype}
  \UseMicrotypeSet[protrusion]{basicmath} % disable protrusion for tt fonts
}{}
\makeatletter
\@ifundefined{KOMAClassName}{% if non-KOMA class
  \IfFileExists{parskip.sty}{%
    \usepackage{parskip}
  }{% else
    \setlength{\parindent}{0pt}
    \setlength{\parskip}{6pt plus 2pt minus 1pt}}
}{% if KOMA class
  \KOMAoptions{parskip=half}}
\makeatother
\usepackage{longtable,booktabs,array}
\usepackage{calc} % for calculating minipage widths
% Correct order of tables after \paragraph or \subparagraph
\usepackage{etoolbox}
\makeatletter
\patchcmd\longtable{\par}{\if@noskipsec\mbox{}\fi\par}{}{}
\makeatother
% Allow footnotes in longtable head/foot
\IfFileExists{footnotehyper.sty}{\usepackage{footnotehyper}}{\usepackage{footnote}}
\makesavenoteenv{longtable}
\usepackage{graphicx}
\makeatletter
\newsavebox\pandoc@box
\newcommand*\pandocbounded[1]{% scales image to fit in text height/width
  \sbox\pandoc@box{#1}%
  \Gscale@div\@tempa{\textheight}{\dimexpr\ht\pandoc@box+\dp\pandoc@box\relax}%
  \Gscale@div\@tempb{\linewidth}{\wd\pandoc@box}%
  \ifdim\@tempb\p@<\@tempa\p@\let\@tempa\@tempb\fi% select the smaller of both
  \ifdim\@tempa\p@<\p@\scalebox{\@tempa}{\usebox\pandoc@box}%
  \else\usebox{\pandoc@box}%
  \fi%
}
% Set default figure placement to htbp
\def\fps@figure{htbp}
\makeatother
% definitions for citeproc citations
\NewDocumentCommand\citeproctext{}{}
\NewDocumentCommand\citeproc{mm}{%
  \begingroup\def\citeproctext{#2}\cite{#1}\endgroup}
\makeatletter
 % allow citations to break across lines
 \let\@cite@ofmt\@firstofone
 % avoid brackets around text for \cite:
 \def\@biblabel#1{}
 \def\@cite#1#2{{#1\if@tempswa , #2\fi}}
\makeatother
\newlength{\cslhangindent}
\setlength{\cslhangindent}{1.5em}
\newlength{\csllabelwidth}
\setlength{\csllabelwidth}{3em}
\newenvironment{CSLReferences}[2] % #1 hanging-indent, #2 entry-spacing
 {\begin{list}{}{%
  \setlength{\itemindent}{0pt}
  \setlength{\leftmargin}{0pt}
  \setlength{\parsep}{0pt}
  % turn on hanging indent if param 1 is 1
  \ifodd #1
   \setlength{\leftmargin}{\cslhangindent}
   \setlength{\itemindent}{-1\cslhangindent}
  \fi
  % set entry spacing
  \setlength{\itemsep}{#2\baselineskip}}}
 {\end{list}}
\usepackage{calc}
\newcommand{\CSLBlock}[1]{\hfill\break\parbox[t]{\linewidth}{\strut\ignorespaces#1\strut}}
\newcommand{\CSLLeftMargin}[1]{\parbox[t]{\csllabelwidth}{\strut#1\strut}}
\newcommand{\CSLRightInline}[1]{\parbox[t]{\linewidth - \csllabelwidth}{\strut#1\strut}}
\newcommand{\CSLIndent}[1]{\hspace{\cslhangindent}#1}
\setlength{\emergencystretch}{3em} % prevent overfull lines
\providecommand{\tightlist}{%
  \setlength{\itemsep}{0pt}\setlength{\parskip}{0pt}}
\usepackage{bookmark}
\IfFileExists{xurl.sty}{\usepackage{xurl}}{} % add URL line breaks if available
\urlstyle{same}
\hypersetup{
  pdftitle={GapHarness: Obligation-First Minimal Harness Synthesis for API-Only LLM Agents},
  pdfauthor={Anonymous Authors},
  hidelinks,
  pdfcreator={LaTeX via pandoc}}

\title{GapHarness: Obligation-First Minimal Harness Synthesis for
API-Only LLM Agents}
\author{Anonymous Authors}
\date{}

\begin{document}
\maketitle

\section{Abstract}\label{abstract}

API-only LLM agents need external runtime support for retrieval,
execution, state, action, control, and verification. Tool routers often
choose tools directly, conflating what a query requires with which
runtime modules happen to exist. GapHarness separates these questions:
it first infers external obligations, then compiles the lowest-cost
declared runtime harness that covers them. On the 1000-task
human-audited GapBench controlled benchmark, GapHarness matches the
oracle minimal harness under gold obligations. With LLM-inferred
obligations, it outperforms direct and router baselines while exposing
calibration tradeoffs. Registry guarding reduces a systematic
sandbox-action unsupported false-positive failure. Anti-circularity
stress tests show that success depends on declared registry affordances
and meaningful obligation labels.

\section{1. Motivation}\label{motivation}

Tool-using LLM agents and reasoning-action systems show that external
actions can improve model behavior (\citeproc{ref-yao2023react}{Yao et
al. 2023}; \citeproc{ref-schick2023toolformer}{Schick et al. 2023};
\citeproc{ref-patil2023gorilla}{Patil et al. 2023};
\citeproc{ref-qin2024toolllm}{Qin et al. 2024}). But tool choice is only
one layer. A query may require observation beyond the prompt,
deterministic execution, durable state, sandbox action, risk control, or
independent verification. These are obligations, not tool names.

GapHarness asks a narrow systems question:

\begin{quote}
Given a user query, an obligation profile, and a declared registry of
modules, can we compile the minimal harness that satisfies the external
obligations required for a warranted answer or action?
\end{quote}

This is not another general agent framework. It is a compiler-style
layer between obligation inference and runtime module selection.
Minimality is relative to a declared ontology, registry, dependency
model, and cost function.

\section{2. Method}\label{method}

GapHarness uses six obligations: Observation, Execution, State, Action,
Control, and Verification. A profiler emits obligations, required
capabilities, output contracts, and unsupported or clarification
conditions. A registry declares modules, affordances, dependencies, and
cost. The compiler performs exact search over registry subsets and
returns the lowest-cost subset that covers the profile. If no subset
covers the profile, it returns unsupported or clarify rather than
hallucinating support.

\begin{figure}
\centering
\includegraphics[width=0.95\linewidth,height=\textheight,keepaspectratio]{paper/figures/figure1_pipeline_print.png}
\caption{GapHarness pipeline. The system separates obligation profiling,
declared registry lookup, exact compilation, sandbox execution, tracing,
and verification.}
\end{figure}

\section{3. Benchmark and Metrics}\label{benchmark-and-metrics}

GapBench-1000 is a controlled factorial benchmark with human-audited
labels for obligations, required capabilities, oracle minimal harnesses,
expected status, and risk metadata. It is designed to isolate harness
compilation, not to replace open-world answer-level benchmarks such as
GAIA (\citeproc{ref-mialon2024gaia}{Mialon et al. 2024}), AgentBench
(\citeproc{ref-liu2024agentbench}{Liu et al. 2024}), Terminal-Bench
(\citeproc{ref-merrill2026terminalbench}{Merrill et al. 2026}),
WildToolBench (\citeproc{ref-yu2026wildtoolbench}{Yu et al. 2026}), or
MCP-Bench (\citeproc{ref-wang2025mcpbench}{Wang et al. 2025}).

We report success, average cost, oracle cost, minimality regret,
over-harnessing, under-harnessing, wrong-harnessing, and redundancy.

\section{4. Main Results}\label{main-results}

Table 1 shows the gold-obligation compiler result. It proves that the
exact compiler matches oracle minimal harnesses under human-audited
obligations. It does not prove open-world profiling.

\begin{longtable}[]{@{}
  >{\raggedright\arraybackslash}p{(\linewidth - 16\tabcolsep) * \real{0.0857}}
  >{\raggedleft\arraybackslash}p{(\linewidth - 16\tabcolsep) * \real{0.1143}}
  >{\raggedleft\arraybackslash}p{(\linewidth - 16\tabcolsep) * \real{0.1143}}
  >{\raggedleft\arraybackslash}p{(\linewidth - 16\tabcolsep) * \real{0.1143}}
  >{\raggedleft\arraybackslash}p{(\linewidth - 16\tabcolsep) * \real{0.1143}}
  >{\raggedleft\arraybackslash}p{(\linewidth - 16\tabcolsep) * \real{0.1143}}
  >{\raggedleft\arraybackslash}p{(\linewidth - 16\tabcolsep) * \real{0.1143}}
  >{\raggedleft\arraybackslash}p{(\linewidth - 16\tabcolsep) * \real{0.1143}}
  >{\raggedleft\arraybackslash}p{(\linewidth - 16\tabcolsep) * \real{0.1143}}@{}}
\toprule\noalign{}
\begin{minipage}[b]{\linewidth}\raggedright
System
\end{minipage} & \begin{minipage}[b]{\linewidth}\raggedleft
N
\end{minipage} & \begin{minipage}[b]{\linewidth}\raggedleft
Success
\end{minipage} & \begin{minipage}[b]{\linewidth}\raggedleft
Cost
\end{minipage} & \begin{minipage}[b]{\linewidth}\raggedleft
Oracle
\end{minipage} & \begin{minipage}[b]{\linewidth}\raggedleft
Regret
\end{minipage} & \begin{minipage}[b]{\linewidth}\raggedleft
Over
\end{minipage} & \begin{minipage}[b]{\linewidth}\raggedleft
Under
\end{minipage} & \begin{minipage}[b]{\linewidth}\raggedleft
Wrong
\end{minipage} \\
\midrule\noalign{}
\endhead
\bottomrule\noalign{}
\endlastfoot
direct & 1000 & 0.20 & 0.00 & 3.67 & -3.67 & 0.00 & 0.74 & 0.00 \\
tool & 1000 & 0.34 & 2.10 & 3.67 & -1.57 & 0.11 & 0.60 & 0.42 \\
difficulty & 1000 & 0.43 & 3.46 & 3.67 & -0.21 & 0.28 & 0.51 & 0.16 \\
always\_full & 1000 & 0.94 & 16.00 & 3.67 & 12.33 & 0.94 & 0.00 &
0.00 \\
GH-gold & 1000 & 1.00 & 3.67 & 3.67 & 0.00 & 0.00 & 0.00 & 0.00 \\
\end{longtable}

\section{5. LLM Profiler and Registry
Guard}\label{llm-profiler-and-registry-guard}

With LLM-inferred obligations on held-out test800, GapHarness reaches
0.89 success at 3.59 average cost, outperforming direct and router
baselines. The main remaining failure is profiler calibration. Phase 2C
adds a deterministic registry guard for a systematic sandbox/mock action
failure where the LLM profiler lowered local sandbox actions into
unsupported real-world side effects.

\begin{longtable}[]{@{}
  >{\raggedright\arraybackslash}p{(\linewidth - 16\tabcolsep) * \real{0.0882}}
  >{\raggedright\arraybackslash}p{(\linewidth - 16\tabcolsep) * \real{0.0882}}
  >{\raggedleft\arraybackslash}p{(\linewidth - 16\tabcolsep) * \real{0.1176}}
  >{\raggedleft\arraybackslash}p{(\linewidth - 16\tabcolsep) * \real{0.1176}}
  >{\raggedleft\arraybackslash}p{(\linewidth - 16\tabcolsep) * \real{0.1176}}
  >{\raggedleft\arraybackslash}p{(\linewidth - 16\tabcolsep) * \real{0.1176}}
  >{\raggedleft\arraybackslash}p{(\linewidth - 16\tabcolsep) * \real{0.1176}}
  >{\raggedleft\arraybackslash}p{(\linewidth - 16\tabcolsep) * \real{0.1176}}
  >{\raggedleft\arraybackslash}p{(\linewidth - 16\tabcolsep) * \real{0.1176}}@{}}
\toprule\noalign{}
\begin{minipage}[b]{\linewidth}\raggedright
Split
\end{minipage} & \begin{minipage}[b]{\linewidth}\raggedright
Profiler
\end{minipage} & \begin{minipage}[b]{\linewidth}\raggedleft
N
\end{minipage} & \begin{minipage}[b]{\linewidth}\raggedleft
Success
\end{minipage} & \begin{minipage}[b]{\linewidth}\raggedleft
Cost
\end{minipage} & \begin{minipage}[b]{\linewidth}\raggedleft
Regret
\end{minipage} & \begin{minipage}[b]{\linewidth}\raggedleft
Over
\end{minipage} & \begin{minipage}[b]{\linewidth}\raggedleft
Under
\end{minipage} & \begin{minipage}[b]{\linewidth}\raggedleft
Unsupported FP
\end{minipage} \\
\midrule\noalign{}
\endhead
\bottomrule\noalign{}
\endlastfoot
test800 & 2B & 800 & 0.89 & 3.59 & -0.09 & 0.14 & 0.09 & 56 \\
test800 & 2C guard & 800 & 0.94 & 3.98 & 0.30 & 0.15 & 0.03 & 12 \\
\end{longtable}

\begin{figure}
\centering
\includegraphics[width=0.9\linewidth,height=\textheight,keepaspectratio]{paper/figures/figure4_registry_guard_unsupported_fp_reduction.png}
\caption{Registry guard correction. Registry guarding reduces
unsupported false positives from 56 to 12 on held-out GapBench test800,
at the cost of higher regret.}
\end{figure}

\section{6. Anti-Circularity Stress
Tests}\label{anti-circularity-stress-tests}

Registry perturbation removes key modules one at a time. For relevant
60-task subsets, base success is 1.00 and perturbed success is 0.00 for
each removed module. The dominant missing capability matches the removed
affordance.

Gold label permutation corrupts 200 supported profiles while the
verifier checks original labels. Correct labels yield 1.00 success.
Permuted labels reduce success to 0.17 and raise under-harnessing to
0.83 and wrong-harnessing to 0.79. This is an anti-circularity stress
test, not a realistic label-noise model.

Negative controls show that pure-language and tool-bait prompts do not
trigger GapHarness over-harnessing. Tool Router and Difficulty Router
over-harness tool-bait prompts at 0.51; Always-full over-harnesses by
construction.

\begin{longtable}[]{@{}
  >{\raggedright\arraybackslash}p{(\linewidth - 6\tabcolsep) * \real{0.2308}}
  >{\raggedright\arraybackslash}p{(\linewidth - 6\tabcolsep) * \real{0.2308}}
  >{\raggedleft\arraybackslash}p{(\linewidth - 6\tabcolsep) * \real{0.3077}}
  >{\raggedright\arraybackslash}p{(\linewidth - 6\tabcolsep) * \real{0.2308}}@{}}
\toprule\noalign{}
\begin{minipage}[b]{\linewidth}\raggedright
Stress test
\end{minipage} & \begin{minipage}[b]{\linewidth}\raggedright
Condition
\end{minipage} & \begin{minipage}[b]{\linewidth}\raggedleft
N
\end{minipage} & \begin{minipage}[b]{\linewidth}\raggedright
Main result
\end{minipage} \\
\midrule\noalign{}
\endhead
\bottomrule\noalign{}
\endlastfoot
Registry perturbation & six module removals & 60 each & success drops
1.00 to 0.00 \\
Gold permutation & corrupted labels & 200 & success drops 1.00 to
0.17 \\
Tool-bait negative control & GapHarness guarded & 100 & over-harness
0.00 \\
Tool-bait negative control & Tool Router & 100 & over-harness 0.51 \\
\end{longtable}

\section{7. Boundaries}\label{boundaries}

GAIA-Transfer gold smoke reaches 1.00 success under gold obligation
labels, but this is obligation-transfer only and not GAIA answer
solving. GAIA registry-guarded reaches 0.56 success with high over- and
under-harnessing, a useful limitation result.
Terminal-Bench-obligation50 is an appendix scaffold with labels pending
audit, not Terminal-Bench solving.

The executor is deterministic sandbox/mock infrastructure. It does not
perform irreversible file edits, real API calls, payments, emails,
deployments, or production changes. GapBench is controlled and
factorial, not a complete real-world benchmark. LLM profiling remains
imperfect.

\section{8. Conclusion}\label{conclusion}

GapHarness reframes agent harnessing as two separable problems:
obligation inference and minimal runtime compilation. Under gold
obligations, the compiler matches oracle minimal harnesses. Under
LLM-inferred obligations, it improves over simple baselines while
exposing calibration tradeoffs. Registry perturbation, label
permutation, and negative controls show that the system is sensitive to
registry affordances and obligation semantics rather than tool keywords
alone.

\section*{References}\label{references}
\addcontentsline{toc}{section}{References}

\protect\phantomsection\label{refs}
\begin{CSLReferences}{1}{0}
\bibitem[\citeproctext]{ref-liu2024agentbench}
Liu, Xiao, Hao Yu, Hanchen Zhang, Yifan Xu, Xuanyu Lei, Hanyu Lai, Yu
Gu, et al. 2024. {``AgentBench: Evaluating LLMs as Agents.''} In
\emph{International Conference on Learning Representations}.
\url{https://openreview.net/forum?id=zAdUB0aCTQ}.

\bibitem[\citeproctext]{ref-merrill2026terminalbench}
Merrill, Mike A., Alexander G. Shaw, Nicholas Carlini, Boxuan Li, Harsh
Raj, Ivan Bercovich, Lin Shi, et al. 2026. {``Terminal-Bench:
Benchmarking Agents on Hard, Realistic Tasks in Command Line
Interfaces.''} \emph{arXiv Preprint arXiv:2601.11868}.
\url{https://arxiv.org/abs/2601.11868}.

\bibitem[\citeproctext]{ref-mialon2024gaia}
Mialon, Grégoire, Clémentine Fourrier, Craig Swift, Thomas Wolf, Yann
LeCun, and Thomas Scialom. 2024. {``GAIA: A Benchmark for General AI
Assistants.''} In \emph{International Conference on Learning
Representations}. \url{https://arxiv.org/abs/2311.12983}.

\bibitem[\citeproctext]{ref-patil2023gorilla}
Patil, Shishir G., Tianjun Zhang, Xin Wang, and Joseph E. Gonzalez.
2023. {``Gorilla: Large Language Model Connected with Massive APIs.''}
\emph{arXiv Preprint arXiv:2305.15334}.
\url{https://arxiv.org/abs/2305.15334}.

\bibitem[\citeproctext]{ref-qin2024toolllm}
Qin, Yujia, Shihao Liang, Yining Ye, Kunlun Zhu, Lan Yan, Yaxi Lu,
Yankai Lin, et al. 2024. {``ToolLLM: Facilitating Large Language Models
to Master 16000+ Real-World APIs.''} In \emph{International Conference
on Learning Representations}.
\url{https://openreview.net/forum?id=dHng2O0Jjr}.

\bibitem[\citeproctext]{ref-schick2023toolformer}
Schick, Timo, Jane Dwivedi-Yu, Roberto Dessì, Roberta Raileanu, Maria
Lomeli, Luke Zettlemoyer, Nicola Cancedda, and Thomas Scialom. 2023.
{``Toolformer: Language Models Can Teach Themselves to Use Tools.''} In
\emph{Advances in Neural Information Processing Systems}.
\url{https://arxiv.org/abs/2302.04761}.

\bibitem[\citeproctext]{ref-wang2025mcpbench}
Wang, Zhenting, Qi Chang, Hemani Patel, Shashank Biju, Cheng-En Wu, Quan
Liu, Aolin Ding, et al. 2025. {``MCP-Bench: Benchmarking Tool-Using LLM
Agents with Complex Real-World Tasks via MCP Servers.''} \emph{arXiv
Preprint arXiv:2508.20453}. \url{https://arxiv.org/abs/2508.20453}.

\bibitem[\citeproctext]{ref-yao2023react}
Yao, Shunyu, Jeffrey Zhao, Dian Yu, Nan Du, Izhak Shafran, Karthik
Narasimhan, and Yuan Cao. 2023. {``ReAct: Synergizing Reasoning and
Acting in Language Models.''} In \emph{International Conference on
Learning Representations}. \url{https://arxiv.org/abs/2210.03629}.

\bibitem[\citeproctext]{ref-yu2026wildtoolbench}
Yu, Peijie, Wei Liu, Yifan Yang, Jinjian Li, Zelong Zhang, Xiao Feng,
and Feng Zhang. 2026. {``Benchmarking LLM Tool-Use in the Wild.''}
\emph{arXiv Preprint arXiv:2604.06185}.
\url{https://arxiv.org/abs/2604.06185}.

\end{CSLReferences}

\end{document}
